\documentclass[a4paper, 12pt]{article}

\usepackage{utils}

\renewcommand*{\today}{10 April 2026}

\begin{document}

\hotbox{Algèbre bilinéaire}{CM 10}{\today}

\section{Groupe orthogonal}

\begin{definition}
    Soit $(E, \langle \rangle)$ un $\K$-espace euclidien vectoriel de dimension finie. Le \textbf{groupe orthogonal} de $E$ est l'ensemble
    $$
        O(E, \langle \rangle) = \{f \in \mathrm{end}(E) \mid \langle f(x), f(y) \rangle = \langle x, y \rangle \text{ pour tous } x, y \in E\}
    $$
\end{definition}

\begin{proposition}{}{}
    Soit $(E, \langle \rangle)$ un $\K$-espace euclidien vectoriel de dimension finie et $B$ une base orthonormée de $E$.
    $$
    ^t Mat_B(f) Mat_B(f) = I_n \iff f \in O(E)
    $$
    ou autrement dit, $M \in \mathcal{M}_n(\K)$ est une matrice orthogonale si et seulement si $^t M M = I_n$.
    
    On sait que $det(M) = \pm 1$ pour une matrice orthogonale $M$.
\end{proposition}

\subsection{Réduction des endomorphismes orthogonaux}

\begin{proposition}{}{}
    Soit $f \in \mathrm{end}(E)$ un endomorphisme orthogonal d'un $\K$-espace euclidien vectoriel de dimension finie.

    Si $f$ possède une valeur propre réelle $\lambda$, alors $\lambda = \pm 1$ et
    $$
    \mathrm{Ker}(f - \lambda id_E) \perp \mathrm{Ker}(f + \lambda id_E)
    $$
\end{proposition}

\begin{corollaire}{}{}
    Si $dim(E) = 1$ alors
    $$
    O^+(E) = \{id_E\} \quad \text{ et } \quad O^-(E) = \{-id_E\}
    $$
\end{corollaire}

\begin{proposition}{}{}
    Soit $(E, \langle \rangle)$ un $\K$-espace euclidien vectoriel de dimension 2.

    Les éléments de $O^+(E)$ s'appellent les rotations vectorielles.

    $O^+(E)$ est un groupe abélien.

    Si $f \in O^+(E)$ et $\mathcal{B}$ est une base orthonormée de $E$, alors
    $$
    Mat_{\mathcal{B}}(f) = \begin{pmatrix}
        \cos \theta & -\sin \theta \\
        \sin \theta & \cos \theta
    \end{pmatrix}
    $$
    pour un certain $\theta \in \R \mod 2\pi$.

    La matrice de $f$ est la même dans toute autre base orthonormée de même orientation que $\mathcal{B}$.

    Les éléments de $O^+(E)$ ne sont pas diagonalisables sauf pour $id_E$ et $-id_E$.
\end{proposition}

\begin{proposition}{}{}
    Soit $(E, \langle \rangle)$ un $\K$-espace euclidien vectoriel de dimension 2.

    Les éléments de $O^-(E)$ s'appellent les symétries vectorielles de $E$.

    $O^-(E)$ n'est pas un groupe.

    Si $f \in O^-(E)$ et $\mathcal{B}$ est une base orthonormée de $E$, alors
    $$
    Mat_{\mathcal{B}}(f) = \begin{pmatrix}
        \cos \theta & \sin \theta \\
        \sin \theta & -\cos \theta
    \end{pmatrix}
    $$
    pour un certain $\theta \in \R \mod 2\pi$.

    La matrice de $f$ est la même dans toute autre base orthonormée de même orientation que $\mathcal{B}$.

    Les éléments de $O^-(E)$ sont diagonalisables.
\end{proposition}

\begin{lemme}{}{}
    Soit $(E, \langle \rangle)$ un $\K$-espace euclidien vectoriel de dimension finie et $V \subset E$ un sous-espace vectoriel de $E$.
    Soit $f \in O(E)$
    $$
    f(V) \subset V \implies f(V^\perp) \subset V^\perp
    $$
\end{lemme}

\begin{lemme}{}{}
    Soit $(E, \langle \rangle)$ un $\K$-espace euclidien vectoriel de dimension finie et $f \in \mathrm{end}(E)$ un endomorphisme de $E$.

    $f$ stabilise un sous-espace vectoriel de dimension 1 ou 2 de $E$.
\end{lemme}

\begin{theoreme}{}{}
    Soit $(E, \langle \rangle)$ un $\K$-espace euclidien vectoriel de dimension finie et $f \in O(E)$ un endomorphisme orthogonal de $E$.

    Il existe une base orthonormée de $E$ dans laquelle la matrice de $f$ est de la forme
    $$
    \begin{pmatrix}
        -I_p & 0 & 0 & \cdots & 0 \\
        0 & I_q & 0 & \cdots & 0 \\
        0 & 0 & R_{\theta_1} & \cdots & 0 \\
        \vdots & \vdots & \ddots & \vdots \\
        0 & 0 & 0 & \cdots & R_{\theta_k} \\
    \end{pmatrix}
    $$
    où $R_{\theta_i}$ est une matrice de rotation de l'espace euclidien de dimension 2, $I_p$ est la matrice identité de l'espace euclidien de dimension $p$ et $n = 2k + p + q$.

    De plus, les entiers $p, q, k$ et les angles $\theta_1, \ldots, \theta_k$ sont uniques.

    En particulier, $f \in O^+(E) \iff p \text{ est pair}$
\end{theoreme}

\end{document}